\subsection{Discussion}
The main takeaway that can be taken from these experiments is that good performance on clean text can hide a dependence on its exact spelling and encoding. Character features generally tolerate the tested changes better than word and BPE features, probably because of their higher "density" but they still struggle when much of the input is modified. DistilBERT also shows how different the outcome can depend on the attack: its preprocessing removes the tested invisible characters, while other substitutions remain damaging. A single average score would hide these differences.

The defense results show how much can be recovered by changing the way the input is handled, without retraining the classifier. This makes simple preprocessing worth considering before moving to a more complex model. The almost complete recovery with full repair needs some caution, though: the repair map covers the characters used by the attacks. I would expect a less complete result with unfamiliar substitutions, and this experiment cannot tell us how safely the same rules would handle legitimate Unicode text in other languages, some possible next steps may be the study of smarter, less destructive and/or more generalized defenses for homoglyphs attacks.

Jev changes very little under these attacks, even when the local models lose much of their performance. The uncertainty is mainly about why it behaves this way and how far that behavior would generalize. Unknown training data, possible exposure to the datasets, and any preprocessing behind the API make it difficult to attribute the result to the model's architecture. Normalization similar to the defenses used here could explain part of its robustness, but I have no evidence that the API applies it. The experiment measures the behavior of the service as a whole.

